\scanpage{34}
\margin{33}
du sommet correspondant de $C^*$. De façon plus précise, les côtés correspondent biunivoquement aux \uline{arêtes} du polygone local des arêtes \deleted{de} issues de $\unsure{\sigma|s}$, ou encore aux sommets de $\Gamma$ incidents à la face.

Plus précisément encore, le polygône comb. associé à cette comp. connexe du bord est can. isom. au \uline{dual} du polygône combinatoire correspondant à la face envisagée de $C$. \unsure{(}

Dans le cas qui nous intéresse, provenant d'un syst. de ps. droites $\Sigma$, en prenant la carte \unsure{structurelle} de toutes les positions relatives, il se trouve que l'ordre de toutes ces faces est (un diviseur de) 4. De plus, parmi les \inserted{quatre} côtés des diverses comp. connexes du bord, correspondant donc à un \uline{coin} de la carte $C$, il y en a un qui est \uline{privilégié}

Mais la propriété de ``type'' qui le distingue ne dépend pas \inserted{(en général)} que du sommet correspondant de $C$ \inserted{i.e. de $\Gamma$}, mais bien du coin \deleted{\unsure{qui la}} particulier envisagé. Donc il n'y a \deleted{\unsure{pas}} guère d'espoir, en général, de réaliser \inserted{le graphe sous-jacent à} $C$ comme réunion de deux \deleted{cartes $C_0,C$} graphes $\gamma,\gamma^*$, correspondants topologiquement à deux \deleted{\unsure{graphes d}} \inserted{cartes} \unsure{duales} sur $\mathcal X$, duales l'une de l'autre.

Soit $\underline\Sigma=(X,\Sigma)$ un syst. de pseudo-droites, et $L\deleted{D}$ une position relative (non triviale). On prend une \deleted{petite}\inserted{étroite} bande \inserted{(de Möbius) $\beta_L$} autour de $L$, dont le bord $\widetilde L$ peut être vu comme rev. à deux feuillets de $\unsure{L|D}$ (\uline{\unsure{rev.|rec.} des vis}), et le complémentaire de $\mathring\beta_L$ est un disque, \deleted{$\unsure{D_L=\Delta_L}$} soit $\Delta_L$. Toute $D_i\in\Sigma$ coupe $\partial\Delta_L=\deleted{\unsure{D}}\widetilde L$ \deleted{\unsure{sur}} en exactement deux points, et sa trace sur \deleted{$\unsure{D_{\Delta_L}}$} sur $\Delta_L$ est un segment (ainsi que sa trace dans $\beta_L$).
